\section{Voice Interaction}
\label{sec:voice}
\begin{figure}[!t]
\centering
\begin{tikzpicture}[
  font=\scriptsize, node distance=3.4mm,
  b/.style={draw, rounded corners=2pt, align=center, text width=6.6cm, minimum height=5.5mm, inner sep=2.5pt, fill=black!3},
  pf/.style={-{Latex[length=1.7mm]}}
]
\node[b] (press) {mic control: press-and-hold, or quick tap $\to$ hands-free toggle};
\node[b, below=of press] (pre) {preflight: \code{acquireMicStream()} --- try OS default, else enumerate every \code{audioinput} and try each with \code{deviceId:\{ideal\}}};
\node[b, below=of pre] (rec) {\code{SpeechRecognition} (continuous, interim results); silence timers $7$\,s / $3$\,s; hard cap $25$\,s};
\node[b, below=of rec] (tx) {final transcript accumulated, submitted once};
\node[b, below=of tx] (send) {written into the input field; auto-submit after $350$\,ms via the \emph{same} chat send path};
\node[b, below=of send] (rag) {RAG assistant (Section~\ref{sec:rag}); answer optionally read aloud via \code{speechSynthesis}};
\draw[pf] (press)--(pre); \draw[pf] (pre)--(rec); \draw[pf] (rec)--(tx);
\draw[pf] (tx)--(send); \draw[pf] (send)--(rag);
\end{tikzpicture}
\caption{Voice-to-RAG interaction pipeline. Speech-to-text is
browser-native and never contacts the back end directly; text-to-speech is
optional and browser-native. Voice-navigation command parsing is not
implemented.}
\label{fig:voice}
\end{figure}


Voice input is provided entirely in the browser through the Web Speech API,
consistent with the constraint of incurring no server-side voice cost
(Fig.~\ref{fig:voice}). There is \emph{zero} voice code in the back end.

\subsection{Speech-to-Text}
The \code{useSpeechRecognition} hook wraps \code{SpeechRecognition} /
\code{webkitSpeechRecognition}. Feature detection runs post-mount to avoid a
server/client hydration mismatch. Recognition runs in continuous mode with
interim results so a brief mid-sentence pause does not end the session
early; final fragments are accumulated and submitted exactly once, when
recognition ends (user release, a $3$-second post-speech silence timeout, a
$7$-second initial-silence timeout, or a $25$-second hard cap). Error codes
are mapped to actionable messages distinguishing permission denial, no
speech, no microphone, and an engine that cannot open the OS default input.

A microphone-acquisition routine, \code{acquireMicStream}, first attempts
an unconstrained \code{getUserMedia(\{audio:true\})}; on
\code{NotFoundError}, \code{OverconstrainedError}, or \code{NotReadableError}
it enumerates every real \code{audioinput} device and tries each with
\code{deviceId:\{ideal\}}. This specifically addresses the case where the
operating system's default communication device is stale (e.g.\ a
disconnected Bluetooth headset) while the built-in microphone is present
and usable as a non-default device. It never retries on a permission
denial. A development-only diagnostic panel reports device counts before
and after permission, track state, and which acquisition step succeeded.

On a successful transcript, the text is written into the chat input and
auto-submitted after $350$\,ms through the \emph{same} send path as typed
input --- voice never talks to the back end directly. Recognition language
follows the selected UI language.

\subsection{Text-to-Speech}
An optional ``read answers aloud'' toggle uses \code{speechSynthesis}. The
hook ranks available voices (preferring a regional and then a
network/neural voice, never requiring one), strips markdown and citation
artifacts before speaking, and splits the answer into moderate-length
utterances for natural pacing. This affects only spoken output; the
displayed text is untouched.

\subsection{What Is Not Implemented}
There is no voice-navigation command parsing: a transcript such as ``take
me to the library'' is passed to the chat pipeline as an ordinary query and
is \emph{not} intercepted as a navigation intent. There are no server-side
speech endpoints. These omissions are consistent with the browser-only
design intent; they are narrower in scope than an aspirational ``voice
navigation'' feature.
